Nilpotency: $P^2(a) = P\bigl((a\cdot\et)\er\bigr) = (a\cdot\et)(\er\cdot\et)\,\er = 0$, since $\er \cdot \et = 0$. The group law is bilinearity plus $P^2 = 0$: $(1 + sP)(1 + s'P) = 1 + (s + s')P + ss'P^2$. Inversion is the case $s' = -s$. The metric: $\hb^{-1}(\et) = \et + \beta(\et\cdot\et)\er = \et + \beta\er$, and $\hb^{-1}(\er) = \er$ since $\er \cdot \et = 0$; then $g(\et,\et) = (\et + \beta\er)^2 = \et^2 + \beta^2\er^2 = 1 - \beta^2$ (the cross terms cancel by anticommutation of orthogonal vectors, \claimref{std-algebra.orthogonal-anticommute}), $g(\et,\er) = (\et + \beta\er)\cdot\er = \beta\,\er^2 = -\beta$, and $g(\er,\er) = -1$; the angular generators are untouched by $P$. For (4): the folding computation is short enough to repeat. The null directions of $g$ in the $(t,r)$ plane are the $a = \et + \lambda\er$ with $0 = g(a,a) = (\et + (\lambda + \beta)\er)^{2} = 1 - (\lambda + \beta)^{2}$, that is $g(a,a) = (1-\beta^2) - 2\lambda\beta - \lambda^2$, whose roots are $\lambda = -\beta \pm 1$. As $r \to 0$, $\beta \to \infty$ and both normalised null directions $\tfrac{1}{-\beta\pm1}\et + \er \to \er$: the cone collapses onto the single direction $\er$, the kernel direction, null for the rescaled boundary form in which the collapse is measured \cite[Prop.~2.6]{ButlerT5} (\claimref{singularity.degeneration-exactly}) (for $g$ itself $\er$ has norm $-1$ at every $r > 0$). That the resulting obstruction is type N and forces Rule 1 (\claimref{singularity.type-n-only-type-n}), and that no gauge choice removes it (\claimref{singularity.non-removability-invariant-argument}), is T5's Steps 1-2 \cite{ButlerT5}, quoted here rather than reproved.
